% ============================================================
% FILE: latex_additions.tex
%
% Integration instructions:
%   - Section A: Insert at the beginning of Chapter 3 (Architecture Concepts),
%     immediately after the \chapter{Architecture Concepts} heading and before
%     Table 3.1 (Revision History).
%   - Section B: Replace the placeholder body of Chapter 6 (Test And Debug),
%     i.e. replace everything after Table 6.1 (Revision History).
%   - Add the bibliography entries in Section C to your references.bib file.
%
% Both sections reference the existing specification content and assume the
% acronyms/glossary already defined in the document (CPU, ALU, etc.).
% ============================================================


% ============================================================
% SECTION A — Chapter 3: Brief Introduction (insert before Table 3.1)
% ============================================================

The architecture of the RWU64IM microcontroller is designed around a set of explicit
structural decisions that directly govern both its functional behaviour and its security
posture.  The processor implements the RV64I base integer instruction set, extended
with the RV64M, RV64A, and RV64D instruction subsets, and is built around a scalar,
in-order pipeline operating at \SI{80}{\mega\hertz}.  All memory accesses use physical
addressing without virtual memory translation; the system exposes a single machine-mode
privilege level with no user or supervisor separation.

A key consequence of this architectural minimalism is the elimination of entire classes
of microarchitectural vulnerabilities.  Transient execution attacks such as
Meltdown~\cite{lipp2018meltdown} and Spectre~\cite{kocher2019spectre} are predicated
on the simultaneous presence of speculative or out-of-order execution, a hardware-enforced
isolation boundary, and a shared observable side-channel (typically the cache hierarchy).
The absence of privilege separation and virtual memory in this design structurally
removes the preconditions required for Meltdown and for Spectre Variant 2 (Branch
Target Injection).  Spectre Variant 1 (Bounds Check Bypass) remains conditionally
relevant in multi-component software deployments that share the physical address space
without hardware-enforced boundaries; a cache timing side-channel is present due to the
shared data cache.  The security implications of these properties are analysed in
detail in Chapter~\ref{chap:testdebug}.

The architectural concepts described in this chapter — the CPU and FPU pipeline
structure, bus topology, memory organisation, interrupt model, clock strategy, and
register layout — are presented with this combined functional and security context
in mind.


% ============================================================
% SECTION B — Chapter 6: Test And Debug (full replacement body)
% ============================================================
%
% Replace the placeholder text "This Section will be completed later. ..."
% with the following content. Keep Table 6.1 (Revision History) in place.
% ============================================================

\label{chap:testdebug}

This chapter analyses the security properties of the RWU64IM microcontroller with
respect to transient execution attacks, documents the conditions under which the
architecture is and is not vulnerable, and specifies both the functional verification
strategy and the security test approach.

% ----------------------------------------------------------
\section{Security Analysis: Transient Execution Vulnerabilities}
% ----------------------------------------------------------

Transient execution attacks exploit the gap between what a processor commits
architecturally and what it executes microarchitecturally during speculative or
out-of-order processing.  Two published attack families — Meltdown
(CVE-2017-5754)~\cite{lipp2018meltdown} and Spectre
(CVE-2017-5753, CVE-2017-5715)~\cite{kocher2019spectre} — represent the primary
threat model evaluated against this design.

% ----------------------------------------------------------
\subsection{Attack Preconditions}
% ----------------------------------------------------------

Three architectural conditions must be simultaneously present for a transient
execution attack to be instantiable~\cite{kocher2019spectre,lipp2018meltdown}:

\begin{enumerate}
  \item \textbf{Speculative or out-of-order (OoO) execution.}  Instructions must be
        capable of executing before their architectural validity has been confirmed.
        Without this, no transient execution window exists.
  \item \textbf{Shared, observable cache.}  The cache must be accessible to both the
        speculating execution path and the attacker's measurement code, enabling
        timing-based inference of transiently accessed data~\cite{gruss2016flushflush}.
  \item \textbf{Hardware-enforced isolation boundary (attack-class specific).}
        Meltdown additionally requires a privilege boundary and a page-fault mechanism
        to create a deferred exception window.  Spectre Variant~2 additionally requires
        separate execution contexts sharing a branch predictor state.
\end{enumerate}

Vulnerability arises from the composition of these conditions; the removal of any
single required precondition is sufficient to invalidate the corresponding attack
class~\cite{intel2018analysis}.

% ----------------------------------------------------------
\subsection{Meltdown Analysis}
% ----------------------------------------------------------

Meltdown exploits the delay between speculative execution of a faulting memory load
and the enforcement of the access permission that causes the fault.  During this
transient window the illegally loaded data may propagate into microarchitectural state
(cache), from which it is exfiltrated via a timing side-channel~\cite{lipp2018meltdown}.

The mechanism requires:
\begin{itemize}
  \item A hardware-enforced privilege boundary (user/kernel or equivalent),
  \item Virtual memory with page tables and per-page permission bits,
  \item A page-fault exception that is raised \emph{after} the dependent speculative
        instructions have already modified cache state, and
  \item A shared cache as the covert exfiltration channel.
\end{itemize}

The RWU64IM design implements none of the first three conditions.  The processor
operates exclusively in machine mode with physical addressing; there are no page
tables, no MMU, no \texttt{satp} register, and no page-fault exceptions.  Every
memory access to the physically addressed space succeeds architecturally; no faulting
load can be issued, and therefore no transient fault window is created.

\textbf{Result:} Meltdown is structurally impossible on this design.  The flat memory
model means that any address can be read directly — this is a consequence of the
architectural choice, not a mitigated vulnerability.

% ----------------------------------------------------------
\subsection{Spectre Analysis}
% ----------------------------------------------------------

Spectre exploits branch misprediction to cause the processor to speculatively execute
unintended code paths.  The speculatively accessed data is encoded into cache state
and subsequently recovered by timing measurement~\cite{kocher2019spectre}.

\subsubsection{Variant 1 — Bounds Check Bypass (CVE-2017-5753)}

Variant~1 requires a predictable conditional branch (e.g.\ a bounds check), a branch
predictor that can be trained by an attacker, a cache-encodable out-of-bounds access
on the mispredicted path, and an isolation boundary between attacker and victim to
make the leakage meaningful.

The RWU64IM processor includes a hardware branch predictor and an in-order-issued,
out-of-order-executing pipeline with a Reorder Buffer, satisfying the execution and
prediction prerequisites.  A shared data cache is present, providing the covert
channel.  However, the system exposes no hardware-enforced isolation boundary:
all software executes in the same machine-mode trust domain with uniform physical
memory access.  In a single-application deployment there is therefore no meaningful
security boundary to cross.

\textbf{Qualification:} If the system runs multiple co-resident software components
(e.g.\ a cooperative real-time kernel with multiple tasks sharing the address space),
a software-defined boundary emerges.  In this scenario Spectre~Variant~1 is feasible
because the hardware provides no mechanism to prevent one task from speculatively
reading another task's data.  The standard mitigation is insertion of speculation
barriers (\texttt{fence} instructions in RISC-V) at security-sensitive branch sites,
or replacement of conditional bounds checks with constant-time index-masking
(\texttt{x = x \& (size-1)}).

\subsubsection{Variant 2 — Branch Target Injection (CVE-2017-5715)}

Variant~2 poisons the Branch Target Buffer (BTB) of a victim execution context to
redirect speculative indirect-branch execution to an attacker-chosen gadget.  This
requires separate, isolated execution contexts sharing predictor hardware state.

The RWU64IM design implements a single machine-mode privilege level.  There is no
separate victim context and no isolation boundary between predictor-training code and
victim code.  BTB poisoning is technically possible but has no meaningful target: all
code already executes in the same domain and can access all memory directly.

\textbf{Result:} Spectre Variant~2 is not applicable.

% ----------------------------------------------------------
\subsection{Cache Timing Side-Channel}
% ----------------------------------------------------------

A cache timing side-channel exists independently of the Meltdown and Spectre attack
classes~\cite{gruss2016flushflush,arm2018cache}.  It requires only:
\begin{itemize}
  \item A shared cache whose fill state depends on access patterns,
  \item A measurable timing difference between a cache hit and a cache miss, and
  \item A mechanism for the attacker to trigger or observe data-dependent memory
        accesses.
\end{itemize}

The data cache of the RWU64IM satisfies all three conditions.  The primary
measurement primitive on RISC-V is the \texttt{rdcycle} CSR instruction, which is
accessible in machine mode.  Typical observed thresholds are: L1 cache hit~$\leq 15$
cycles; DRAM access~$\geq 150$ cycles~\cite{kocher2019spectre}.  The Flush+Reload
technique~\cite{lipp2018meltdown} is realisable using cache eviction loops (the
RISC-V Zicbom extension \texttt{cbo.inval}~\cite{riscv2022zicbom} provides an
explicit invalidation instruction when available).

\textbf{Result:} A cache timing side-channel is present and functional.  It does not
by itself constitute a Meltdown or Spectre exploit on this design, but represents a
residual microarchitectural leakage mechanism that must be addressed when co-resident
software components with differing trust levels are deployed.

% ----------------------------------------------------------
\subsection{Vulnerability Summary}
% ----------------------------------------------------------

\begin{table}[h]
\centering
\caption{Transient Execution Vulnerability Assessment}
\label{tab:vuln_summary}
\begin{tabular}{llp{7cm}}
\hline
\textbf{Vulnerability} & \textbf{Status} & \textbf{Reason} \\
\hline
Meltdown (CVE-2017-5754)     & Not applicable        & No privilege boundary; no page-fault mechanism; physical addressing only \\
Spectre V1 (CVE-2017-5753)   & Conditionally possible & Hardware preconditions (branch predictor, cache) present; requires software-defined isolation boundary to be exploitable \\
Spectre V2 (CVE-2017-5715)   & Not applicable        & No separate execution context; no victim isolation domain \\
Cache timing side-channel     & Present               & Shared cache with measurable hit/miss latency difference \\
\hline
\end{tabular}
\end{table}

% ----------------------------------------------------------
\section{Countermeasures}
% ----------------------------------------------------------

\subsection{Software-Level Mitigations}

\paragraph{Speculation barriers.}
The RISC-V \texttt{fence} instruction serialises memory operations and prevents
speculative execution beyond the barrier point.  It should be inserted immediately
after bounds-check branches at software isolation boundaries:

\begin{verbatim}
if (x < size) {
    asm volatile ("fence");
    val = array[x];
}
\end{verbatim}

\paragraph{Index masking (branch elimination).}
Bounds checks can be replaced with arithmetic masking for power-of-two array sizes,
eliminating the conditional branch and thus the branch predictor dependency:

\begin{verbatim}
x   = x & (size - 1);
val = array[x];
\end{verbatim}

\paragraph{Constant-time access patterns.}
Security-sensitive code should access data unconditionally and select results via
arithmetic masks rather than conditional branches, removing data-dependent cache
footprints.  This is standard practice in cryptographic implementations.

\paragraph{Cycle-counter restriction.}
Clearing the \texttt{CY} bit in the \texttt{mcounteren} CSR prevents unprivileged
code from reading \texttt{rdcycle}, degrading the timing resolution available to
an attacker.  This increases noise in the side-channel but does not eliminate it.

\subsection{Hardware-Level Mitigations}

\paragraph{Cache partitioning.}
The 2-way set-associative cache can be partitioned between software components
(one cache way per component) to prevent cross-component cache-set eviction
observation.  This requires modification to the cache controller and a software
convention for way allocation.

\paragraph{Speculative cache-fill suppression.}
The strongest hardware countermeasure is to prevent cache-line fills from instructions
that are subsequently flushed from the Reorder Buffer (i.e.\ instructions that are
never architecturally committed).  This aligns microarchitectural cache state with
architectural commit semantics, eliminating the persistent side-channel that all
transient execution attacks depend on.  Implementation requires tagging speculative
load operations in the ROB and suppressing the corresponding cache-fill signals until
commit.

\begin{table}[h]
\centering
\caption{Countermeasure Applicability Summary}
\label{tab:countermeasures}
\begin{tabular}{llll}
\hline
\textbf{Countermeasure} & \textbf{Target} & \textbf{Type} & \textbf{Effectiveness} \\
\hline
\texttt{fence} instruction barriers  & Spectre V1   & Software & High \\
Index masking                         & Spectre V1   & Software & High \\
Constant-time access patterns         & Side-channel & Software & High \\
Cycle-counter restriction (\texttt{mcounteren}) & Side-channel & Software & Medium \\
Cache way partitioning                & Side-channel & Hardware & High \\
Speculative cache-fill suppression    & All classes  & Hardware & Very High \\
\hline
\end{tabular}
\end{table}

% ----------------------------------------------------------
\section{Functional Verification Strategy}
% ----------------------------------------------------------

The functional test and debug infrastructure is built around the toolchain described
in Chapter~5.  Verification covers three levels: unit-level RTL simulation, firmware
integration testing, and regression testing.

\subsection{RTL Simulation}

Hardware behaviour is verified using \textbf{Xilinx Vivado XSIM} with SystemVerilog
testbenches.  Simulation is invoked either in console mode (for automation and CI
pipelines) or in GUI mode with waveform capture for interactive debugging:

\begin{verbatim}
# Console mode
xsim <snapshot> -runall

# GUI mode with waveform configuration
xsim <snapshot> --gui --tclbatch sim/config/xsim_cfg2.tcl
\end{verbatim}

Each functional block — CPU core, FPU, interrupt controller, UART, QSPI, GPIO,
Wishbone interconnect — has a dedicated testbench exercising its interface signals
and expected register behaviour.

\subsection{FPU Verification}

FPU correctness is verified against the IEEE~754-2008 standard using the
\textbf{Berkeley TestFloat} suite~\cite{hauser2010testfloat}, which provides
exhaustive corner-case vectors covering all five rounding modes, subnormal operands,
NaN propagation (quiet and signalling), and all five exception flags (\texttt{NX},
\texttt{UF}, \texttt{OF}, \texttt{DZ}, \texttt{NV}).  Execution latencies
(FADD:~3~cycles, FMUL:~4~cycles, FDIV:~20--24~cycles, FSQRT:~25--30~cycles) are
verified by cycle-accurate simulation.

\subsection{Firmware Integration Tests}

Application-level tests are compiled using the RISC-V GCC toolchain configured for
\texttt{-march=rv64imd} and loaded into instruction memory via the JTAG interface.
Each test covers a specific instruction class or peripheral interaction and produces
a pass/fail result observable on the simulation waveform or a UART output channel.

\subsection{Regression Testing}

The CTest framework integrated into the CMake build system aggregates all unit,
integration, firmware, and IP-level tests into a single parallel regression run:

\begin{verbatim}
ctest --test-dir build -j
\end{verbatim}

All test results are reported per-target, enabling automated detection of regressions
across pipeline stages.

% ----------------------------------------------------------
\section{Debug Interface}
% ----------------------------------------------------------

On-chip debug access is provided through the JTAG interface (IEEE~1149.1-compliant
pins: \texttt{JTAG\_TCK}, \texttt{JTAG\_TMS}, \texttt{JTAG\_TDI}, \texttt{JTAG\_TDO},
\texttt{JTAG\_TRST}).  The JTAG TAP controller allows:

\begin{itemize}
  \item Direct loading of binary images into I-MEM,
  \item Read access to internal CPU registers and CSRs for debug inspection, and
  \item System reset assertion via the TAP reset path.
\end{itemize}

A placeholder for a formal JTAG register-level description and TAP state diagram
is reserved for the final document revision:

\begin{figure}[h]
\centering
% TODO: Insert JTAG TAP state diagram and register map here.
\fbox{\parbox{0.85\textwidth}{\centering\vspace{2cm}
\textit{[Placeholder: JTAG TAP State Diagram]}
\vspace{2cm}}}
\caption{JTAG TAP controller state diagram (to be completed).}
\label{fig:jtag_tap}
\end{figure}

\begin{figure}[h]
\centering
% TODO: Insert pipeline simulation waveform showing fetch-decode-execute-writeback
%       timing for a representative instruction sequence.
\fbox{\parbox{0.85\textwidth}{\centering\vspace{2cm}
\textit{[Placeholder: Pipeline Simulation Waveform]}
\vspace{2cm}}}
\caption{Representative pipeline simulation waveform (XSIM output).}
\label{fig:pipeline_wave}
\end{figure}


% ============================================================
% SECTION C — Bibliography entries (add to references.bib)
% ============================================================
%
% @inproceedings{lipp2018meltdown,
%   author    = {Lipp, Moritz and Schwarz, Michael and Gruss, Daniel and
%                Prescher, Thomas and Haas, Werner and Horn, Jann and
%                Mangard, Stefan and Kocher, Paul and Genkin, Daniel and
%                Yarom, Yuval and Hamburg, Mike},
%   title     = {{Meltdown: Reading Kernel Memory from User Space}},
%   booktitle = {27th USENIX Security Symposium},
%   year      = {2018},
%   pages     = {973--990},
%   publisher = {USENIX Association}
% }
%
% @inproceedings{kocher2019spectre,
%   author    = {Kocher, Paul and Horn, Jann and Fogh, Anders and Genkin, Daniel and
%                Gruss, Daniel and Haas, Werner and Hamburg, Mike and Lipp, Moritz and
%                Mangard, Stefan and Prescher, Thomas and Schwarz, Michael and
%                Yarom, Yuval},
%   title     = {{Spectre Attacks: Exploiting Speculative Execution}},
%   booktitle = {40th IEEE Symposium on Security and Privacy},
%   year      = {2019},
%   pages     = {1--19},
%   publisher = {IEEE}
% }
%
% @inproceedings{gruss2016flushflush,
%   author    = {Gruss, Daniel and Maurice, Cl{\'e}mentine and
%                Fogh, Anders and Mangard, Stefan},
%   title     = {{Flush+Flush: A Fast and Stealthy Cache Attack}},
%   booktitle = {Detection of Intrusions and Malware, and Vulnerability
%                Assessment (DIMVA)},
%   year      = {2016},
%   pages     = {279--299},
%   publisher = {Springer}
% }
%
% @techreport{intel2018analysis,
%   author      = {{Intel Corporation}},
%   title       = {{Analysis of Speculative Execution Side Channels}},
%   institution = {Intel Corporation},
%   year        = {2018},
%   note        = {White Paper, Revision~4.0}
% }
%
% @techreport{arm2018cache,
%   author      = {{ARM Limited}},
%   title       = {{Cache Speculation Side-channels}},
%   institution = {ARM Limited},
%   year        = {2018},
%   note        = {Whitepaper, Version~2.4}
% }
%
% @techreport{riscv2022zicbom,
%   author      = {{RISC-V International}},
%   title       = {{RISC-V Cache Management Operation ISA Extensions (Zicbom)}},
%   institution = {RISC-V International},
%   year        = {2022},
%   note        = {Version~1.0}
% }
%
% @techreport{hauser2010testfloat,
%   author      = {Hauser, John R.},
%   title       = {{TestFloat: Floating-Point Verification}},
%   institution = {UC Berkeley},
%   year        = {2010},
%   note        = {\url{http://www.jhauser.us/arithmetic/TestFloat.html}}
% }
% ============================================================
